El mismo circuito, excitado con un escalón de \SI{1}{\volt}. La
Figura~\ref{fig:escalon} lleva cuatro curvas y usa la convención de color de
Xtal en los dos ejes que tiene: \textbf{el color dice qué señal es} (ámbar la
entrada, verde la salida) y \textbf{el trazo dice de dónde salió el dato}
(sólida la teórica, con marcadores la simulada, punteada la medida). Por eso las
dos salidas comparten el verde sin confundirse.

Se leen tres cosas de un vistazo:

\begin{itemize}[itemsep=2pt, topsep=4pt]
  \item \textbf{Sobrepico.} La salida llega a \SI{1.41}{\volt} contra un escalón
        de \SI{1}{\volt}: \SI{41}{\percent}, contra el \SI{45}{\percent} que
        predice el modelo ideal. La diferencia es, otra vez, $R_L$.
  \item \textbf{Oscilación amortiguada.} El período entre máximos es de
        \SI{0.96}{\milli\second}, o sea \SI{1.04}{\kilo\hertz}: es $\omega_d$,
        apenas por debajo de $\fnat$, como manda
        $\omega_d = \sqrt{\omega_0^2 - \alpha^2}$.
  \item \textbf{Tiempo de establecimiento.} La envolvente cae como
        $e^{-\alpha t}$ con $\alpha = \SI{1650}{\per\second}$, así que entrar en
        la banda del \SI{2}{\percent} lleva $4/\alpha \approx
        \SI{2.4}{\milli\second}$. Coincide con lo medido.
  \end{itemize}

\noindent
La curva medida arranca antes de $t = 0$: son las muestras de \emph{pretrigger}
del osciloscopio. El disparo está en el flanco de CH1, que es lo que fija el
origen de tiempos común a las tres fuentes.
